% Weekly report — <repo>, <week>
%
% Seeded from the toolkit's template, then filled in. The structure is fixed:
% title, overview, one frame per theme, what is next. A theme is a strand of
% work a reader would name — never one bullet per commit, which turns the deck
% back into the log it was supposed to summarise.
%
% Nothing here needs a package beyond a stock TeX install, so the file compiles
% wherever it is opened. A sibling preamble.tex, if this repo has one, is
% pulled in below: that is where a project's own theme, colours, fonts and
% title-page branding belong, so this file can stay generated.
\documentclass[aspectratio=169]{beamer}
\usepackage[T1]{fontenc}
\usepackage{lmodern}
\usetheme{default}
\setbeamertemplate{navigation symbols}{}
\IfFileExists{preamble.tex}{\input{preamble.tex}}{}

\title{<repo>}
\subtitle{Weekly report --- <week>}
\date{<period>}

\begin{document}

\begin{frame}
  \titlepage
\end{frame}

\begin{frame}{Overview}
  % Three to five lines, in the language of the goal rather than of the commit
  % log: what moved, and how far. Someone who reads only this frame should
  % know where the project stands.
  \begin{itemize}
    \item
  \end{itemize}
\end{frame}

% ---------------------------------------------------------------------------
% One frame per theme, three to five in total. Copy the frame per theme, give
% each a title naming the strand of work, and delete the spare.
%
% The footnotesize line at the bottom is the evidence trail — pull request
% numbers, paths, run identifiers — for the reader who wants to follow a claim
% back to where it happened. Drop the line rather than leave it empty.
% ---------------------------------------------------------------------------
\begin{frame}{Theme}
  \begin{itemize}
    \item
  \end{itemize}
  \vfill
  {\footnotesize }
\end{frame}

\begin{frame}{Next}
  % What the coming week takes on. Short — this frame is a pointer, not a plan.
  \begin{itemize}
    \item
  \end{itemize}
\end{frame}

\end{document}
